\documentclass[12pt,a4paper]{article}

% Packages
\usepackage[utf8]{inputenc}
\usepackage[T1]{fontenc}
\usepackage{geometry}
\usepackage{graphicx}
\usepackage{amsmath}
\usepackage{amssymb}
\usepackage{booktabs}
\usepackage{hyperref}
\usepackage{fancyhdr}
\usepackage{titlesec}
\usepackage{enumitem}
\usepackage{xcolor}
\usepackage{listings}
\usepackage{float}
\usepackage{caption}
\usepackage{subcaption}

\geometry{margin=2.5cm, top=3cm, bottom=2.5cm}

% Header with logo
\pagestyle{fancy}
\fancyhf{}
\fancyhead[L]{
\includegraphics[height=1.2cm]{uorlogo.jpg}}
\fancyhead[R]{\small Explainable AI --- MTech Data Science}
\fancyfoot[C]{\thepage}
\renewcommand{\headrulewidth}{0.4pt}
\renewcommand{\footrulewidth}{0.4pt}

% Code listing style
\lstset{
    basicstyle=\ttfamily\small,
    breaklines=true,
    frame=single,
    backgroundcolor=\color{gray!10},
    keywordstyle=\color{blue},
    commentstyle=\color{green!60!black},
    stringstyle=\color{red!70!black},
    language=Python
}

\begin{document}

% Title Page
\begin{titlepage}
    \centering
    
\includegraphics[width=4cm]{uorlogo.jpg}\\[1cm]
    {\Large \textbf{Nirma University}}\\[0.3cm]
    {\large Institute of Technology}\\[0.3cm]
    {\large Department of Computer Science \& Engineering}\\[1.5cm]
    {\LARGE \textbf{Explainable AI}}\\[0.5cm]
    {\Large \textbf{Lab 1: Explainability in Intrinsically Explainable Models}}\\[1.5cm]
    {\large \textbf{(Regression, Decision Trees, Random Forest + SHAP)}}\\[2cm]
    \begin{tabular}{ll}
        \textbf{Student Name:} & Soham Dave \\[0.3cm]
        \textbf{Roll Number:} & 25MCD005 \\[0.3cm]
        \textbf{Program:} & MTech Data Science, Semester 2 \\[0.3cm]
        \textbf{Subject:} & Explainable AI \\[0.3cm]
        \textbf{Subject Guide:} & Dr. Nileshkumar Patel \\
    \end{tabular}
    \vfill
    {\large \today}
\end{titlepage}

\tableofcontents
\newpage

% -------------------------------------------------------
\section{Objective}
The objective of this lab is to demonstrate explainability in \textbf{intrinsically interpretable} machine learning models. Three models are explored:
\begin{enumerate}
    \item \textbf{Multiple Linear Regression} --- interpreting coefficients as feature importance
    \item \textbf{Decision Tree Classifier} --- visualizing tree structure and Gini importance
    \item \textbf{Random Forest Classifier with SHAP} --- using SHapley Additive exPlanations for post-hoc explainability
\end{enumerate}

% -------------------------------------------------------
\section{Datasets}

\subsection{Cellphone Dataset}
\begin{itemize}
    \item \textbf{Task:} Regression (predicting cellphone price)
    \item \textbf{Samples:} 161 entries
    \item \textbf{Features (12):} Sale, weight, resolution, ppi, cpu core, cpu freq, internal mem, ram, RearCam, Front\_Cam, battery, thickness
    \item \textbf{Target:} Price (integer)
    \item \textbf{Preprocessing:} Standard scaling (zero mean, unit variance) applied to all features
\end{itemize}

\subsection{MBA Admissions Dataset}
\begin{itemize}
    \item \textbf{Task:} Classification (predicting admission decision)
    \item \textbf{Samples:} 1,000 (with valid admission labels out of 6,194 total)
    \item \textbf{Features (34 after encoding):} gpa, gmat, work\_exp, gender, international, major, race, work\_industry (one-hot encoded)
    \item \textbf{Target:} Admission decision (categorical)
    \item \textbf{Preprocessing:} Rows with missing \texttt{admission} dropped; categorical features one-hot encoded
\end{itemize}

\subsection{Wine Quality Dataset}
\begin{itemize}
    \item \textbf{Task:} Multi-class classification (predicting wine quality score)
    \item \textbf{Samples:} 1,599 entries
    \item \textbf{Features (11):} Fixed acidity, volatile acidity, citric acid, residual sugar, chlorides, free/total sulfur dioxide, density, pH, sulphates, alcohol
    \item \textbf{Target:} Quality (integer score)
    \item \textbf{Split:} 80/20 train-test split
\end{itemize}

% -------------------------------------------------------
\section{Methodology}

\subsection{Part 1: Multiple Linear Regression}

\subsubsection{Model Description}
Multiple Linear Regression models the target variable as a weighted linear combination of input features:
\begin{equation}
    \hat{y} = \beta_0 + \beta_1 x_1 + \beta_2 x_2 + \cdots + \beta_p x_p
\end{equation}

where $\beta_0$ is the intercept and $\beta_1, \ldots, \beta_p$ are the regression coefficients. Each coefficient $\beta_j$ quantifies the expected change in the target for a one-unit change in feature $x_j$, holding all other features constant.

\subsubsection{Implementation Steps}
\begin{enumerate}
    \item Loaded \texttt{Cellphone.csv} and inspected dataset structure using \texttt{df.info()}
    \item Separated the target variable (\texttt{Price}) from features; dropped \texttt{Product\_id}
    \item Applied \texttt{StandardScaler} to normalize all features to zero mean and unit variance
    \item Fitted \texttt{LinearRegression} from scikit-learn on the entire dataset
    \item Extracted model coefficients (\texttt{model.coef\_}) for explainability analysis
\end{enumerate}

\subsubsection{Explainability Approach}
Since features are standardized, the \textbf{absolute magnitude of coefficients} serves as a direct measure of feature importance. A larger absolute coefficient indicates a stronger influence on the predicted price.

\begin{figure}[H]
    \centering
    \includegraphics[width=0.9\textwidth]{images_1/output_1.png}
    \caption{Feature importance derived from Linear Regression coefficients (absolute values). Internal memory, battery, and thickness show the highest importance.}
    \label{fig:lr_importance}
\end{figure}

\subsubsection{Results Interpretation}
Key coefficients obtained:
\begin{itemize}
    \item \texttt{internal mem}: $+179.47$ --- strongest positive predictor of price
    \item \texttt{battery}: $+171.56$ --- second strongest positive predictor
    \item \texttt{thickness}: $-161.06$ --- strong negative effect (thinner phones cost more)
    \item \texttt{ram}: $+149.57$ and \texttt{ppi}: $+135.18$ --- significant positive contributors
    \item \texttt{Sale}: $-35.07$ and \texttt{weight}: $-35.15$ --- minor negative effects
\end{itemize}

The coefficient signs and magnitudes provide a complete, transparent explanation of the model's pricing logic.

% -------------------------------------------------------
\subsection{Part 2: Decision Tree Classifier}

\subsubsection{Model Description}
A Decision Tree recursively partitions the feature space using axis-aligned splits that maximize class purity. At each internal node, the algorithm selects the feature and threshold that best separates the classes (using Gini impurity in this case):
\begin{equation}
    \text{Gini}(t) = 1 - \sum_{k=1}^{K} p_k^2
\end{equation}
where $p_k$ is the proportion of class $k$ at node $t$.

\subsubsection{Implementation Steps}
\begin{enumerate}
    \item Loaded \texttt{MBA.csv} and filtered to 1,000 rows with valid admission labels
    \item One-hot encoded categorical features (gender, international, major, race, work\_industry), yielding 34 features
    \item Trained \texttt{DecisionTreeClassifier} with \texttt{random\_state=42}
    \item Extracted \texttt{feature\_importances\_} (mean decrease in Gini impurity)
    \item Visualized the complete tree structure using \texttt{plot\_tree}
\end{enumerate}

\subsubsection{Explainability Approach}
Decision trees are inherently interpretable because:
\begin{itemize}
    \item Every prediction can be traced through a sequence of if-then rules
    \item Feature importance is quantified by the total reduction in Gini impurity contributed by each feature across all splits
    \item The tree structure itself serves as the explanation
\end{itemize}

\begin{figure}[H]
    \centering
    \includegraphics[width=\textwidth]{images_1/output_2.png}
    \caption{Full Decision Tree visualization for MBA admission prediction. Nodes are colored by majority class, with split conditions and sample counts shown.}
    \label{fig:dt_tree}
\end{figure}

\begin{figure}[H]
    \centering
    \includegraphics[width=0.85\textwidth]{images_1/output_3.png}
    \caption{Feature importance scores from Decision Tree. GPA, GMAT, and work experience dominate the decision process.}
    \label{fig:dt_importance}
\end{figure}

\subsubsection{Results Interpretation}
\begin{itemize}
    \item \textbf{GPA (0.200)} is the single most important feature for admission
    \item \textbf{GMAT (0.175)} and \textbf{work experience (0.118)} follow closely
    \item Demographic features (gender, race, international status) have minimal importance ($<0.015$), suggesting the model relies primarily on academic and professional merit
\end{itemize}

% -------------------------------------------------------
\subsection{Part 3: Random Forest with SHAP}

\subsubsection{Model Description}
A Random Forest is an ensemble of $B$ decision trees, each trained on a bootstrap sample with random feature subsets. The final prediction is the majority vote:
\begin{equation}
    \hat{y} = \text{mode}\{h_1(\mathbf{x}), h_2(\mathbf{x}), \ldots, h_B(\mathbf{x})\}
\end{equation}

While individual trees are interpretable, the ensemble is a \textbf{black box}. SHAP (SHapley Additive exPlanations) provides post-hoc explainability.

\subsubsection{SHAP Theory}
SHAP values are based on Shapley values from cooperative game theory. For a prediction $f(\mathbf{x})$:
\begin{equation}
    f(\mathbf{x}) = \phi_0 + \sum_{j=1}^{p} \phi_j
\end{equation}
where $\phi_0$ is the base value (average prediction) and $\phi_j$ is the SHAP value for feature $j$, representing its marginal contribution.

\subsubsection{Implementation Steps}
\begin{enumerate}
    \item Loaded \texttt{winequality-red.csv} and split into train/test (80/20)
    \item Trained \texttt{RandomForestClassifier} with default hyperparameters
    \item Initialized \texttt{shap.TreeExplainer} for efficient SHAP computation on tree ensembles
    \item Computed SHAP values for all test samples
    \item Generated a SHAP summary plot
\end{enumerate}

\begin{figure}[H]
    \centering
    \includegraphics[width=0.9\textwidth]{images_1/output_4.png}
    \caption{SHAP summary plot for Random Forest on wine quality. Each row is a feature; each dot is a test sample. Color indicates feature value (red=high, blue=low); x-axis shows SHAP value (impact on prediction).}
    \label{fig:shap_summary}
\end{figure}

\subsubsection{Results Interpretation}
\begin{itemize}
    \item \textbf{Alcohol} has the widest spread of SHAP values, making it the most influential feature for wine quality
    \item \textbf{Volatile acidity} and \textbf{sulphates} are the next most important
    \item High alcohol content (red dots) consistently pushes predictions toward higher quality (positive SHAP values)
    \item High volatile acidity (red dots) pushes predictions toward lower quality (negative SHAP values)
\end{itemize}

% -------------------------------------------------------
\section{Conclusion}
This lab demonstrated that intrinsically explainable models offer transparent and interpretable predictions:
\begin{itemize}
    \item \textbf{Linear Regression} provides direct coefficient-based feature importance
    \item \textbf{Decision Trees} offer rule-based explanations through their tree structure
    \item \textbf{SHAP} bridges the gap between complex models (Random Forest) and interpretability by providing theoretically grounded feature attributions
\end{itemize}

These methods form the foundation of Explainable AI, enabling practitioners to understand, trust, and debug their machine learning models.

\end{document}
